\subsection{Adversarial Alternating Optimization}
\label{sec:adv_alternating_training}

Let $\mathbf{x}_i \in \mathbb{R}^{N \times C}$ denote the $i$-th point cloud in a mini-batch $\mathcal{B}$, and let
$f_{\theta}(\cdot)$ be the place descriptor network parameterized by $\theta$. The adversarial occlusion generator
$G_{\phi}(\cdot)$, parameterized by $\phi$, predicts a set of vehicle-like cuboids and applies a structured occlusion
transform to the input point cloud. The overall objective is formulated as a min-max game:
\begin{equation}
\min_{\theta} \max_{\phi} \; \mathcal{L}_{\mathrm{place}}\!\left(f_{\theta}(T_{\phi}(\mathbf{x}))\right),
\label{eq:minmax_objective}
\end{equation}
where $T_{\phi}(\cdot)$ denotes the occlusion transform induced by the generator. In practice, we optimize this
objective by alternating between a generator update step and a descriptor update step in every training iteration.

\paragraph{Batch-hard place recognition loss.}
For each anchor $i$ in the mini-batch, let $\mathcal{P}(i)$ and $\mathcal{N}(i)$ denote the positive and negative
index sets provided by the training tuple construction. Given embeddings $\mathbf{z}_i = f_{\theta}(\mathbf{x}_i)$,
we define the Euclidean distance
\begin{equation}
d_{ij} = \left\lVert \mathbf{z}_i - \mathbf{z}_j \right\rVert_2.
\end{equation}
The batch-hard triplet loss is then written as
\begin{equation}
\mathcal{L}_{\mathrm{place}}
=
\frac{1}{|\mathcal{V}|}
\sum_{i \in \mathcal{V}}
\left[
\max_{p \in \mathcal{P}(i)} d_{ip}
-
\min_{n \in \mathcal{N}(i)} d_{in}
+
\delta
\right]_+,
\label{eq:batch_hard_triplet}
\end{equation}
where $\delta$ is the margin, $[\cdot]_+ = \max(\cdot, 0)$, and $\mathcal{V}$ is the set of valid anchors that have
at least one positive and one negative sample within the current mini-batch.

\paragraph{Generator parameterization and occlusion transform.}
For each input point cloud, the generator predicts $M$ cuboids
\begin{equation}
G_{\phi}(\mathbf{x}_i)
\rightarrow
\left\{
\left(\mathbf{c}_{i,m}, \mathbf{s}_{i,m}, \psi_{i,m}\right)
\right\}_{m=1}^{M},
\end{equation}
where $\mathbf{c}_{i,m} \in \mathbb{R}^3$, $\mathbf{s}_{i,m} \in \mathbb{R}^3$, and $\psi_{i,m}$ denote the center,
size, and yaw angle of the $m$-th cuboid, respectively. During training, an active-box count $k_i \in \{1,\dots,M\}$
is randomly sampled for each sample, and exactly $k_i$ cuboids are activated to produce the occlusion pattern.

The generator outputs both a differentiable soft mask and a hard binary mask. We therefore distinguish two occlusion
operators:
\begin{equation}
\tilde{\mathbf{x}}^{\,s}_i = T_{\phi}^{\mathrm{soft}}(\mathbf{x}_i),
\qquad
\tilde{\mathbf{x}}^{\,h}_i = T_{\phi}^{\mathrm{hard}}(\mathbf{x}_i),
\end{equation}
where the soft transform is used for generator optimization, while the hard transform is used for descriptor
optimization. The hard transform may additionally insert synthetic points on the surfaces of active cuboids to mimic
dynamic object returns.

\paragraph{Generator objective.}
In the generator update step, the descriptor network is frozen and only $\phi$ is optimized. The generator aims to
maximize the place recognition difficulty of the occluded point clouds while remaining geometrically plausible. The
generator loss is defined as
\begin{equation}
\mathcal{L}_{G}
=
-\mathcal{L}_{\mathrm{place}}
\left(
\left\{
f_{\theta}(\tilde{\mathbf{x}}^{\,s}_i)
\right\}_{i \in \mathcal{B}}
\right)
+
\lambda_{\mathrm{size}} \mathcal{L}_{\mathrm{size}}
+
\lambda_{\mathrm{height}} \mathcal{L}_{\mathrm{height}}
+
\lambda_{\mathrm{range}} \mathcal{L}_{\mathrm{range}}.
\label{eq:generator_loss}
\end{equation}
The negative sign makes the generator adversarial: minimizing $\mathcal{L}_{G}$ is equivalent to maximizing the place
loss of the descriptor network on the occluded inputs.

The regularization terms enforce simple geometric priors on the generated cuboids. Specifically,
\begin{equation}
\mathcal{L}_{\mathrm{size}}
=
\begin{cases}
\mathrm{SmoothL1}\!\left(\mathbf{s}_{i,m}, \mathbf{s}_{\mathrm{car}}\right), & \text{if box sizes are learnable}, \\
0, & \text{if fixed box sizes are used},
\end{cases}
\label{eq:size_prior}
\end{equation}
where $\mathbf{s}_{\mathrm{car}}$ is a nominal vehicle size prior. The height prior is
\begin{equation}
\mathcal{L}_{\mathrm{height}}
=
\frac{1}{BM}
\sum_{i=1}^{B}
\sum_{m=1}^{M}
\left(c^{(z)}_{i,m} + 1.0\right)^2,
\label{eq:height_prior}
\end{equation}
which encourages the cuboid centers to stay close to the drivable vertical band. The range prior is
\begin{equation}
\mathcal{L}_{\mathrm{range}}
=
\frac{1}{BM}
\sum_{i=1}^{B}
\sum_{m=1}^{M}
\left[
\max\!\left(
0,
\sqrt{\left(c^{(x)}_{i,m}\right)^2 + \left(c^{(y)}_{i,m}\right)^2} - 45.0
\right)
\right]^2,
\label{eq:range_prior}
\end{equation}
which discourages the generator from placing cuboids at unrealistically large radial distances.

\paragraph{Descriptor objective.}
In the descriptor update step, the generator is frozen and the same sampled box-count budget is reused to produce a
hard occluded point cloud for each sample. The descriptor is optimized jointly on clean and adversarial inputs:
\begin{equation}
\mathbf{z}^{\,c}_i = f_{\theta}(\mathbf{x}_i),
\qquad
\mathbf{z}^{\,a}_i = f_{\theta}(\tilde{\mathbf{x}}^{\,h}_i).
\end{equation}
The descriptor loss is defined as
\begin{equation}
\mathcal{L}_{F}
=
\mathcal{L}_{\mathrm{clean}}
+
\lambda_{\mathrm{adv}} \mathcal{L}_{\mathrm{adv}}
+
\lambda_{\mathrm{cons}} \mathcal{L}_{\mathrm{cons}},
\label{eq:descriptor_loss}
\end{equation}
where
\begin{align}
\mathcal{L}_{\mathrm{clean}}
&=
\mathcal{L}_{\mathrm{place}}
\left(
\left\{
f_{\theta}(\mathbf{x}_i)
\right\}_{i \in \mathcal{B}}
\right),
\\
\mathcal{L}_{\mathrm{adv}}
&=
\mathcal{L}_{\mathrm{place}}
\left(
\left\{
f_{\theta}(\tilde{\mathbf{x}}^{\,h}_i)
\right\}_{i \in \mathcal{B}}
\right),
\\
\mathcal{L}_{\mathrm{cons}}
&=
\frac{1}{B}
\sum_{i=1}^{B}
\left(
1 -
\frac{
\left\langle \mathbf{z}^{\,c}_i, \mathbf{z}^{\,a}_i \right\rangle
}{
\left\lVert \mathbf{z}^{\,c}_i \right\rVert_2
\left\lVert \mathbf{z}^{\,a}_i \right\rVert_2
}
\right).
\label{eq:consistency_loss}
\end{align}
Here, $\mathcal{L}_{\mathrm{clean}}$ preserves the nominal place-recognition capability on uncorrupted point clouds,
$\mathcal{L}_{\mathrm{adv}}$ explicitly improves robustness under generated occlusion, and
$\mathcal{L}_{\mathrm{cons}}$ constrains the clean and adversarial embeddings of the same scene to remain close in the
descriptor space.

\paragraph{Alternating optimization.}
For each mini-batch, we perform the following two sub-steps:
\begin{align}
\phi &\leftarrow \phi - \eta_{G} \nabla_{\phi} \mathcal{L}_{G},
\\
\theta &\leftarrow \theta - \eta_{F} \nabla_{\theta} \mathcal{L}_{F},
\end{align}
where $\eta_{G}$ and $\eta_{F}$ are the learning rates of the generator and the descriptor network, respectively.
This alternating strategy stabilizes the adversarial game: the generator continuously searches for geometrically
plausible hard cases, while the descriptor network learns to preserve place-level discriminability under both clean and
occluded observations.
